% GENERATED FILE. Edit canonical lessons, then run npm run generate:books.
% canonical-source-hash: fnv1a64:4af724ac6065b3b3
% canonical-lessons: SA-C412-viramah, SA-C412-nivartate, SA-C412-upagacchati, SA-C412-atigacchati, SA-C412-krosah

\chapter{Halts and Turning Back}
\label{ch:sa-halts-and-turning-back}
\hlchaptermodality{412}

\begin{chapteropening}
\textbf{By the end of this chapter:} \emph{I can say where we halted, turned back or went past.}

Complete the last lesson of chapter 412: I can say where we halted, turned back or went past.
\end{chapteropening}

\section[\texorpdfstring{\sk{विरामः} (virāmaḥ)}{विरामः (virāmaḥ)}]{\sk{विरामः} (virāmaḥ) — a halt; a stop}

\label{lesson:SA-C412-viramah}



\begin{quote}
Before the new one: say the Sanskrit for a blockage, then the Sanskrit for keeps away from.
\end{quote}

\subsection*{You'll want to know: \sk{विरामः}}
\textbf{\sk{विरामः}} — \emph{virāmaḥ} — \textquotedblleft{}a halt; a stop\textquotedblright{}.

\begin{quote}
\sk{लोकयानस्य} \sk{विरामः} \sk{ग्रामे} \sk{अस्ति।}
\end{quote}

\begin{quote}
\emph{lokayānasya virāmaḥ grāme asti.}
\end{quote}

\begin{quote}
The bus has a halt at the village.
\end{quote}

A masculine noun.

\begin{cousinweb}
\textbf{\sk{विरामः}} (\emph{virāmaḥ}) is built from \emph{vi}, \textquotedblleft{}away\textquotedblright{}, and the root √ram, \textquotedblleft{}to rest, to delight\textquotedblright{}, which also gives \emph{ramate}, \textquotedblleft{}enjoys\textquotedblright{}; MW gives \emph{virāma} as cessation, end, a stop or pause; a bus halt is that sense in modern use (MW).
\end{cousinweb}

\begin{grammarlens}[title={What you've built}]
One more word for getting around.
\end{grammarlens}

\subsection*{Your turn}
\begin{itemize}
\raggedright
  \item \emph{Say it:} \emph{virāmaḥ}
  \item \emph{Say it:} \emph{virāmaḥ}, once more
  \item \emph{Say it:} \emph{lokayānasya virāmaḥ grāme asti.} — The bus has a halt at the village.
  \item \emph{Recall:} say the Sanskrit for a blockage, then the Sanskrit for keeps away from, then say \emph{virāmaḥ} again
\end{itemize}

\begin{tcolorbox}[breakable,colback=teal!4,colframe=teal!35!black,title={Before you move on}]
What does \sk{विरामः} mean? (\textquotedblleft{}A halt; a stop\textquotedblright{}.) Say it once more.
\end{tcolorbox}
\section[\texorpdfstring{\sk{निवर्तते} (nivartate)}{निवर्तते (nivartate)}]{\sk{निवर्तते} (nivartate) — turns round; turns back}

\label{lesson:SA-C412-nivartate}



\begin{quote}
Before the new one: say the Sanskrit for keeps away from, then the Sanskrit for a halt.
\end{quote}

\subsection*{You'll want to know: \sk{निवर्तते}}
\textbf{\sk{निवर्तते}} — \emph{nivartate} — \textquotedblleft{}turns round; turns back\textquotedblright{}.

\begin{quote}
\sk{वर्षायां} \sk{पथिकः} \sk{मार्गात्} \sk{निवर्तते।}
\end{quote}

\begin{quote}
\emph{varṣāyāṃ pathikaḥ mārgāt nivartate.}
\end{quote}

\begin{quote}
In the rain the wayfarer turns back from the road.
\end{quote}

Past (he …): \sk{सः} \sk{मार्गात्} \sk{निवृत्तवान्}, \emph{saḥ mārgāt nivṛttavān}.

\begin{cousinweb}
\textbf{\sk{निवर्तते}} (\emph{nivartate}) is built from \emph{ni}, \textquotedblleft{}back\textquotedblright{}, and the root √vṛt, \textquotedblleft{}to turn\textquotedblright{}, which also gives \emph{vartate}, \textquotedblleft{}is present\textquotedblright{}; MW gives \emph{ni-vṛt} as to turn back, with the participle \emph{nivṛtta}, \textquotedblleft{}turned back, returned\textquotedblright{} (MW).
\end{cousinweb}

\begin{grammarlens}[title={What you've built}]
One more word for getting around.
\end{grammarlens}

\subsection*{Your turn}
\begin{itemize}
\raggedright
  \item \emph{Say it:} \emph{nivartate}
  \item \emph{Say it:} \emph{nivartate}, once more
  \item \emph{Say it:} \emph{varṣāyāṃ pathikaḥ mārgāt nivartate.} — In the rain the wayfarer turns back from the road.
  \item \emph{Recall:} say the Sanskrit for keeps away from, then the Sanskrit for a halt, then say \emph{nivartate} again
\end{itemize}

\begin{tcolorbox}[breakable,colback=teal!4,colframe=teal!35!black,title={Before you move on}]
What does \sk{निवर्तते} mean? (\textquotedblleft{}Turns round; turns back\textquotedblright{}.) Say it once more.
\end{tcolorbox}
\section[\texorpdfstring{\sk{उपगच्छति} (upagacchati)}{उपगच्छति (upagacchati)}]{\sk{उपगच्छति} (upagacchati) — approaches; comes near}

\label{lesson:SA-C412-upagacchati}



\begin{quote}
Before the new one: say the Sanskrit for a halt, then the Sanskrit for turns round.
\end{quote}

\subsection*{You'll want to know: \sk{उपगच्छति}}
\textbf{\sk{उपगच्छति}} — \emph{upagacchati} — \textquotedblleft{}approaches; comes near\textquotedblright{}.

\begin{quote}
\sk{धूमयानं} \sk{स्थानकम्} \sk{उपगच्छति।}
\end{quote}

\begin{quote}
\emph{dhūmayānaṃ sthānakam upagacchati.}
\end{quote}

\begin{quote}
The train is coming up to the station.
\end{quote}

Past (it …): \sk{धूमयानं} \sk{स्थानकम्} \sk{उपगतवत्}, \emph{dhūmayānaṃ sthānakam upagatavat}.

\begin{cousinweb}
\textbf{\sk{उपगच्छति}} (\emph{upagacchati}) is built from \emph{upa}, \textquotedblleft{}towards, near\textquotedblright{}, and the root √gam, \textquotedblleft{}to go\textquotedblright{}, which also gives \emph{gacchati}, \textquotedblleft{}goes\textquotedblright{}; MW gives \emph{upa-gam} as to go near, approach, arrive at (MW).
\end{cousinweb}

\begin{grammarlens}[title={What you've built}]
One more word for getting around.
\end{grammarlens}

\subsection*{Your turn}
\begin{itemize}
\raggedright
  \item \emph{Say it:} \emph{upagacchati}
  \item \emph{Say it:} \emph{upagacchati}, once more
  \item \emph{Say it:} \emph{dhūmayānaṃ sthānakam upagacchati.} — The train is coming up to the station.
  \item \emph{Recall:} say the Sanskrit for a halt, then the Sanskrit for turns round, then say \emph{upagacchati} again
\end{itemize}

\begin{tcolorbox}[breakable,colback=teal!4,colframe=teal!35!black,title={Before you move on}]
What does \sk{उपगच्छति} mean? (\textquotedblleft{}Approaches; comes near\textquotedblright{}.) Say it once more.
\end{tcolorbox}
\section[\texorpdfstring{\sk{अतिगच्छति} (atigacchati)}{अतिगच्छति (atigacchati)}]{\sk{अतिगच्छति} (atigacchati) — passes by; overtakes}

\label{lesson:SA-C412-atigacchati}



\begin{quote}
Before the new one: say the Sanskrit for turns round, then the Sanskrit for approaches.
\end{quote}

\subsection*{You'll want to know: \sk{अतिगच्छति}}
\textbf{\sk{अतिगच्छति}} — \emph{atigacchati} — \textquotedblleft{}passes by; overtakes\textquotedblright{}.

\begin{quote}
\sk{कारयानं} \sk{लोकयानम्} \sk{अतिगच्छति।}
\end{quote}

\begin{quote}
\emph{kārayānaṃ lokayānam atigacchati.}
\end{quote}

\begin{quote}
The car passes the bus.
\end{quote}

Past (it …): \sk{कारयानं} \sk{लोकयानम्} \sk{अतिगतवत्}, \emph{kārayānaṃ lokayānam atigatavat}.

\begin{cousinweb}
\textbf{\sk{अतिगच्छति}} (\emph{atigacchati}) is built from \emph{ati}, \textquotedblleft{}beyond, past\textquotedblright{}, and the root √gam, \textquotedblleft{}to go\textquotedblright{}; MW gives \emph{ati-gam} as to pass by or over, with the participle \emph{atigata}, \textquotedblleft{}having passed\textquotedblright{} (MW).
\end{cousinweb}

\begin{grammarlens}[title={What you've built}]
One more word for getting around.
\end{grammarlens}

\subsection*{Your turn}
\begin{itemize}
\raggedright
  \item \emph{Say it:} \emph{atigacchati}
  \item \emph{Say it:} \emph{atigacchati}, once more
  \item \emph{Say it:} \emph{kārayānaṃ lokayānam atigacchati.} — The car passes the bus.
  \item \emph{Recall:} say the Sanskrit for turns round, then the Sanskrit for approaches, then say \emph{atigacchati} again
\end{itemize}

\begin{tcolorbox}[breakable,colback=teal!4,colframe=teal!35!black,title={Before you move on}]
What does \sk{अतिगच्छति} mean? (\textquotedblleft{}Passes by; overtakes\textquotedblright{}.) Say it once more.
\end{tcolorbox}
\section[\texorpdfstring{\sk{क्रोशः} (krośaḥ)}{क्रोशः (krośaḥ)}]{\sk{क्रोशः} (krośaḥ) — a kos (a measure of distance)}

\label{lesson:SA-C412-krosah}



\begin{quote}
Before the new one: say the Sanskrit for approaches, then the Sanskrit for passes by.
\end{quote}

\subsection*{You'll want to know: \sk{क्रोशः}}
\textbf{\sk{क्रोशः}} — \emph{krośaḥ} — \textquotedblleft{}a kos (a measure of distance)\textquotedblright{}.

\begin{quote}
\sk{ग्रामः} \sk{इतः} \sk{क्रोशं} \sk{दूरम्} \sk{अस्ति।}
\end{quote}

\begin{quote}
\emph{grāmaḥ itaḥ krośaṃ dūram asti.}
\end{quote}

\begin{quote}
The village is a kos away from here.
\end{quote}

A masculine noun.

\begin{cousinweb}
\textbf{\sk{क्रोशः}} (\emph{krośaḥ}) comes from the root √kruś, \textquotedblleft{}to cry out\textquotedblright{}: MW gives \emph{krośa} as a cry and as a measure of distance, the range of the voice in calling, the kos; it lives on as Hindi \emph{kos}, a distance of some two miles (MW; CDIAL 3611).
\end{cousinweb}

\begin{grammarlens}[title={What you've built}]
One more word for getting around.
\end{grammarlens}

\subsection*{Your turn}
\begin{itemize}
\raggedright
  \item \emph{Say it:} \emph{krośaḥ}
  \item \emph{Say it:} \emph{krośaḥ}, once more
  \item \emph{Say it:} \emph{grāmaḥ itaḥ krośaṃ dūram asti.} — The village is a kos away from here.
  \item \emph{Recall:} say the Sanskrit for approaches, then the Sanskrit for passes by, then say \emph{krośaḥ} again
\end{itemize}

\begin{tcolorbox}[breakable,colback=teal!4,colframe=teal!35!black,title={Before you move on}]
What does \sk{क्रोशः} mean? (\textquotedblleft{}A kos (a measure of distance)\textquotedblright{}.) Say it once more.
\end{tcolorbox}
